\documentclass[11pt,a4paper]{article}

\usepackage[margin=2.5cm]{geometry}
\usepackage{amsmath,amssymb}
\usepackage{booktabs}
\usepackage[round]{natbib}
\usepackage[colorlinks=true,linkcolor=blue,citecolor=blue,urlcolor=blue]{hyperref}

\newcommand{\code}[1]{\texttt{#1}}
\newcommand{\pvap}{p_{\rm vap}}
\newcommand{\Pcond}{P_{\rm cond}}
\newcommand{\MH}{[\mathrm{M/H}]}
\newcommand{\chem}[1]{\mathrm{#1}}

\title{Vapour pressures and condensation curves in \textsc{petitRADTRANS}}
\author{}
\date{\today}

\begin{document}
\maketitle

\section{Overview}

This note lists the saturation vapour pressure curves implemented in
\code{petitRADTRANS/chemistry/vapour\_pressures.py}, together with the
original reference for each. They are used in
\code{chemistry/clouds.py} to locate the base of equilibrium-condensation
clouds.

\paragraph{Notation.}
$T$ is the temperature in K, $\pvap$ the saturation vapour pressure, $P$
the total gas pressure, $\MH$ the metallicity relative to solar in $\log_{10}$
units, and $\bar{\mu}$ the mean molar mass of the atmosphere. The equations
below are written in the form used in the code: some use the natural
logarithm and some $\log_{10}$, and they do not all return the same pressure
unit. The unit of each function's return value is listed in
Table~\ref{tab:summary} and stated with each equation.

\paragraph{From vapour pressure to condensation curve.}
For a condensate $c$ with mass fraction $X_c$ and molar mass $M_c$, the
partial pressure of the condensing gas is $P X_c\bar{\mu}/M_c$. The
condensation curve is the total pressure at which this partial pressure
equals the saturation vapour pressure:
\begin{equation}
  \Pcond(T) = \frac{\pvap(T)}{X_c\,\bar{\mu}/M_c} .
  \label{eq:pcond}
\end{equation}
$X_c$ is either derived from the elemental abundances (set by $\MH$ and
C/O, limited by the least abundant element of the condensate), or given as
a free parameter. The cloud base is placed at the deepest intersection
of the temperature profile with the condensation temperature
$T_{\rm cond}(P)$, which is obtained by inverting Eq.~\eqref{eq:pcond}
numerically on a temperature grid. The intersection is located by linear
interpolation in $\log P$.

The exception is \chem{Mg_2SiO_4} with metallicity dependence
(Sect.~\ref{sec:mg2sio4}), for which the fit already gives $\Pcond$ directly.

Several fits are expressed in terms of the vapour pressure of the
limiting \emph{gas-phase} species (e.g.\ Na, Mn or Zn atoms), rather
than of a hypothetical gas of the condensate molecule. These cases are
noted below.

\begin{table}[h]
\centering
\caption{Summary of implemented curves. ``Unit'' is the pressure unit of
the value returned by the function in \code{vapour\_pressures.py}.}
\label{tab:summary}
\begin{tabular}{lllll}
\toprule
Species & Function & Unit & $\MH$ & Original reference \\
\midrule
Fe(s), Fe(l)        & \code{vapour\_pressure\_fe\_solid/\_liquid} & bar & no  & \citet{barshay1976} \\
\chem{MgSiO_3}      & \code{vapour\_pressure\_mgsio3}      & bar & no  & \citet{barshay1976} \\
\chem{Mg_2SiO_4}    & \code{vapour\_pressure\_mg2sio4}     & bar & yes & \citet{visscher2010} \\
\chem{Na_2S}        & \code{vapour\_pressure\_na2s}        & bar & yes & \citet{visscher2006} \\
KCl                 & \code{vapour\_pressure\_kcl}         & bar & no  & \citet{morley2012} \\
SiO                 & \code{vapour\_pressure\_sio}         & bar & no  & \citet{wetzel2013} \\
\chem{TiO_2}        & \code{vapour\_pressure\_tio2}        & dyn\,cm$^{-2}$ & yes & \citet{marley2021} \\
Cr                  & \code{vapour\_pressure\_cr}          & dyn\,cm$^{-2}$ & no  & \citet{morley2012} \\
ZnS                 & \code{vapour\_pressure\_zns}         & dyn\,cm$^{-2}$ & yes & \citet{visscher2006} \\
MnS                 & \code{vapour\_pressure\_mns}         & dyn\,cm$^{-2}$ & yes & \citet{visscher2006} \\
\chem{NH_3}         & \code{vapour\_pressure\_nh3}         & dyn\,cm$^{-2}$ & no  & \citet{lodders1998} \\
\chem{H_2O}         & \code{vapour\_pressure\_h2o}         & dyn\,cm$^{-2}$ & no  & \citet{buck1981}; \citet{wexler1976,wexler1977} \\
\chem{CH_4}         & \code{vapour\_pressure\_ch4}         & dyn\,cm$^{-2}$ & no  & \citet{lodders1998} \\
\chem{Al_2O_3}      & \code{vapour\_pressure\_al2o3}       & dyn\,cm$^{-2}$ & yes & \citet{morley2024} \\
\chem{CaTiO_3}      & \code{vapour\_pressure\_catio3}      & dyn\,cm$^{-2}$ & yes & \citet{wakeford2017} \\
\chem{CaAl_{12}O_{19}} & \code{vapour\_pressure\_caal12o19} & dyn\,cm$^{-2}$ & yes & \citet{wakeford2017} \\
\chem{SiO_2}        & \code{vapour\_pressure\_sio2}        & dyn\,cm$^{-2}$ & yes & \citet{grant2023} \\
\bottomrule
\end{tabular}
\end{table}

The curves for \chem{TiO_2} and all species below it in
Table~\ref{tab:summary} follow the implementation in the \textsc{virga}
cloud code \citep{batalha2025}, which also compiles the original
references.

\section{Iron and silicates}

\subsection{Fe}
Solid (below the melting point) and liquid iron, from
\citet{barshay1976} as adopted by \citet[][Eqs.~A3a,b]{ackerman2001}:
\begin{align}
  \ln \pvap^{\rm Fe(s)} &= 15.71 - \frac{47664}{T}, \\
  \ln \pvap^{\rm Fe(l)} &= 9.86 - \frac{37120}{T},
\end{align}
with $\pvap$ in bar. \citet{ackerman2001} do not state the unit; the pRT
source notes that, per an erratum, the unit is bar and not
dyn\,cm$^{-2}$. When computing the Fe cloud base, pRT takes the
condensation curve as the smaller of the solid and liquid
$\Pcond$ (Eq.~\ref{eq:pcond}) at each temperature.

\subsection{\texorpdfstring{\chem{MgSiO_3}}{MgSiO3}}
Enstatite, from \citet{barshay1976} as adopted by
\citet[][Eq.~A4]{ackerman2001}:
\begin{equation}
  \ln \pvap^{\rm MgSiO_3} = 25.37 - \frac{58663}{T},
\end{equation}
with $\pvap$ in bar (same erratum as for Fe).

\subsection{\texorpdfstring{\chem{Mg_2SiO_4}}{Mg2SiO4}}
\label{sec:mg2sio4}
If a metallicity is given, the forsterite condensation curve of
\citet{visscher2010} is used,
\begin{equation}
  \frac{10^4}{T_{\rm cond}} = 5.89 - 0.37\log_{10}\Pcond - 0.73\,\MH ,
\end{equation}
with $\Pcond$ in bar. Solved for $\Pcond$, as implemented, this is
\begin{equation}
  \log_{10} \Pcond^{\rm Mg_2SiO_4} = 15.92 - \frac{27027.03}{T} - 1.97\,\MH .
\end{equation}
This is the total pressure of condensation, so Eq.~\eqref{eq:pcond} is not
applied. Here $\MH$ stands for [Fe/H] in \citet{visscher2010}.

If no metallicity is given (the free-abundance case), the function instead
returns the \chem{MgSiO_3} vapour pressure of \citet{barshay1976},
$\ln\pvap = 25.37 - 58663/T$ (bar), which is then used in
Eq.~\eqref{eq:pcond} with the \chem{Mg_2SiO_4} mass fraction and molar mass.

\subsection{SiO}
Solid SiO, from \citet[][Eq.~3]{wetzel2013}, who fitted the laboratory
measurements of \citet{ferguson2008} and their own group:
\begin{equation}
  \ln \pvap^{\rm SiO} = 33.848 - \frac{42648}{T},
\end{equation}
with $\pvap$ in dyn\,cm$^{-2}$ in the original. The code multiplies by
$10^{-6}$ to return bar.

\subsection{\texorpdfstring{\chem{SiO_2}}{SiO2}}
Quartz, as listed in \citet{batalha2025}, citing \citet{grant2023}:
\begin{equation}
  \log_{10} \pvap^{\rm SiO_2} = 13.168 - \frac{28265}{T} - \MH ,
\end{equation}
with $\pvap$ in bar; the code returns dyn\,cm$^{-2}$. This is the
``undepleted gas'' version. The code also mentions a ``depleted gas''
alternative, $13.360 - 28265/T - \MH$, which is not used. The pRT
docstring flags this curve as a placeholder that is not yet suitable for
public use.

\section{Refractory oxides}

\subsection{\texorpdfstring{\chem{TiO_2}}{TiO2}}
From \citet{marley2021}, see also \citet{lodders2002}, as implemented in
\textsc{virga}:
\begin{equation}
  \log_{10} \pvap^{\rm TiO_2} = 13.95 - \frac{38266}{T} - \MH ,
\end{equation}
with $\pvap$ in bar; the code returns dyn\,cm$^{-2}$.

\subsection{\texorpdfstring{\chem{Al_2O_3}}{Al2O3}}
Corundum, from the Sonora Diamondback models \citep{morley2024}, see also
\citet{wakeford2017}:
\begin{equation}
  \log_{10} \pvap^{\rm Al_2O_3} = 15.24 - \frac{41481}{T} - 1.66\,\MH ,
\end{equation}
with $\pvap$ in bar; the code returns dyn\,cm$^{-2}$. The metallicity
coefficient $1.66$ follows the \textsc{virga} source code. Table~2 of
\citet{batalha2025} gives $1.50$. An older fit based on
\citet{wakeford2017}, $\log_{10}\pvap = 17.7 - 45892.6/T - 1.66\,\MH$, is
kept as a comment in the code.

\subsection{\texorpdfstring{\chem{CaTiO_3}}{CaTiO3}}
Perovskite, derived from the condensation curve of \citet{wakeford2017}:
\begin{equation}
  \log_{10} \pvap^{\rm CaTiO_3} = 30.24 - \frac{72160}{T}
  - \log_{10}\!\left(\frac{P}{1\,\mathrm{bar}}\right) - 2\,\MH ,
\end{equation}
with $\pvap$ in bar; the code returns dyn\,cm$^{-2}$ and takes $P$ in
dyn\,cm$^{-2}$. \citet{batalha2025} list the same expression with $-3\,\MH$.
In \code{clouds.py}, $P$ is currently fixed to $1$\,bar, so the pressure
term vanishes.

\subsection{\texorpdfstring{\chem{CaAl_{12}O_{19}}}{CaAl12O19}}
Hibonite, derived from the condensation curve of \citet{wakeford2017}:
\begin{equation}
  \log_{10} \pvap^{\rm CaAl_{12}O_{19}} = 16.46 - \frac{44021}{T}
  - 0.083\log_{10}\!\left(\frac{P}{1\,\mathrm{bar}}\right) - 1.67\,\MH ,
\end{equation}
with $\pvap$ in bar; the code returns dyn\,cm$^{-2}$. As for \chem{CaTiO_3},
$P$ is fixed to $1$\,bar in \code{clouds.py}.

\section{Sulfides, salts and Cr}

The following fits are from \citet{morley2012}, who derived them from the
thermochemical models of \citet{visscher2006} (sulfides) and
\citet{lodders1999} (KCl). Pressures are in bar in the original.

\subsection{Cr}
Vapour pressure of Cr gas over Cr metal \citep[][Eq.~5]{morley2012}:
\begin{equation}
  \log_{10} \pvap^{\rm Cr} = 7.49 - \frac{20592}{T} .
\end{equation}
The code returns dyn\,cm$^{-2}$. The \code{metallicity} argument is
accepted but not used.

\subsection{MnS}
Vapour pressure of Mn gas over condensed MnS
\citep[][Eq.~9, from \citealt{visscher2006}]{morley2012}:
\begin{equation}
  \log_{10} \pvap^{\rm Mn} = 11.5315 - \frac{23810}{T} - \MH .
\end{equation}
The code returns dyn\,cm$^{-2}$.

\subsection{ZnS}
Vapour pressure of Zn gas over condensed ZnS
\citep[][Eq.~15, from \citealt{visscher2006}]{morley2012}:
\begin{equation}
  \log_{10} \pvap^{\rm Zn} = 12.8117 - \frac{15873}{T} - \MH .
\end{equation}
The code returns dyn\,cm$^{-2}$.

\subsection{\texorpdfstring{\chem{Na_2S}}{Na2S}}
\citet[][Eq.~12, from \citealt{visscher2006}]{morley2012} give the partial
pressure of Na over condensed \chem{Na_2S},
$\log_{10} p^{\rm Na} = 8.55 - 13889/T - 0.5\,\MH$ (bar). pRT divides this by
two to obtain the partial pressure of a hypothetical \chem{Na_2S} gas, as in
\citet{charnay2018}. This is valid as long as S is more abundant than Na,
as it is for solar abundance ratios:
\begin{equation}
  \pvap^{\rm Na_2S} = \tfrac{1}{2}\,10^{\,8.55 - 13889/T - 0.5\,\MH} ,
\end{equation}
returned in bar.

In the free-abundance case, the code replaces $\MH$ by a proxy for [Na/H]
derived from the mass fraction. As currently implemented, the mass fraction
cancels out, and the expression evaluated is
\begin{equation}
  \pvap^{\rm Na_2S} = \tfrac{1}{2}\,
  10^{\,8.55 - 13889/T - 0.5\left[\log_{10}(2\bar{\mu}) + 5.7\right]} .
\end{equation}
The code comments suggest that the intended proxy is
$[\mathrm{Na/H}] \approx \log_{10}(2X_{\rm Na_2S}\bar{\mu}/M_{\rm Na_2S}) + 5.7$,
which assumes that the solar abundance is $\log_{10}(\mathrm{Na/H})_\odot \approx -5.7$.

\subsection{KCl}
Vapour pressure of KCl over solid KCl \citep[][Eq.~18]{morley2012}, based
on \citet{lodders1999}:
\begin{equation}
  \log_{10} \pvap^{\rm KCl} = 7.6106 - \frac{11382}{T} ,
\end{equation}
returned in bar.

\section{Ices}

\subsection{\texorpdfstring{\chem{NH_3}}{NH3}}
From \citet{lodders1998}, piecewise around the triple point at 195.4\,K:
\begin{equation}
  \log_{10} \pvap^{\rm NH_3} =
  \begin{cases}
    6.900 - 1588/T, & T < 195.4\,\mathrm{K}\ \text{(solid)},\\
    5.201 - 1248/T, & T \ge 195.4\,\mathrm{K}\ \text{(liquid)},
  \end{cases}
\end{equation}
with $\pvap$ in bar; the code returns dyn\,cm$^{-2}$.

\subsection{\texorpdfstring{\chem{H_2O}}{H2O}}
Three temperature regimes are used, following \citet{ackerman2001}.
By default (\code{do\_buck=True}), the expressions of \citet{buck1981} are
used, with $t \equiv T - 273.16$ the temperature in $^\circ$C and $\pvap$ in
dyn\,cm$^{-2}$:
\begin{equation}
  \pvap^{\rm H_2O} =
  \begin{cases}
    6111.5\,\exp\!\left[\left(23.036 - \dfrac{t}{333.7}\right)\dfrac{t}{t + 279.82}\right],
      & T < 273.16\,\mathrm{K}\ \text{(ice)},\\[2ex]
    6112.1\,\exp\!\left[\left(18.729 - \dfrac{t}{227.3}\right)\dfrac{t}{t + 257.87}\right],
      & 273.16 \le T < 1048\,\mathrm{K}\ \text{(liquid)},\\[2ex]
    6\times10^{8}, & T \ge 1048\,\mathrm{K}.
  \end{cases}
\end{equation}
The constant value above 1048\,K (600\,bar) follows \citet{ackerman2001},
who note that the Buck expressions are unsuitable at higher temperatures.

With \code{do\_buck=False}, the formulations of \citet{wexler1977} for ice
and \citet{wexler1976} for liquid water are used, with coefficients as
given by \citet{flatau1992}. The factor 10 converts Pa to dyn\,cm$^{-2}$:
\begin{align}
  \pvap^{\rm ice} &= 10\,\exp\!\left[\frac{h_0}{T} + h_1 + h_2 T + h_3 T^2
     + h_4 T^3 + h_5 \ln T\right], \\
  \pvap^{\rm liq} &= 10\,\exp\!\left[\frac{g_0}{T^2} + \frac{g_1}{T} + g_2
     + g_3 T + g_4 T^2 + g_5 T^3 + g_6 T^4 + g_7 \ln T\right],
\end{align}
with
\begin{align*}
  (h_0,\dots,h_5) &= (-5.8653696\times10^{3},\ 2.224103300\times10^{1},\
     1.3749042\times10^{-2},\ -3.4031775\times10^{-5},\\
     &\qquad 2.6967687\times10^{-8},\ 6.918651\times10^{-1}),\\
  (g_0,\dots,g_7) &= (-2.9912729\times10^{3},\ -6.0170128\times10^{3},\
     1.887643854\times10^{1},\ -2.8354721\times10^{-2},\\
     &\qquad 1.7838301\times10^{-5},\ -8.4150417\times10^{-10},\
     4.4412543\times10^{-13},\ 2.858487).
\end{align*}
The same temperature ranges and high-temperature cap as above apply.

\subsection{\texorpdfstring{\chem{CH_4}}{CH4}}
From \citet{lodders1998}. The Clausius--Clapeyron equation is integrated
with a temperature-independent difference in heat capacity
$\Delta c_p = c_{p,\rm gas} - c_{p,\rm cond}$ between gas and condensate,
and anchored at the triple point
($T_{\rm tp} = 90.68$\,K, $p_{\rm tp} = 0.11719$\,bar):
\begin{equation}
  \pvap^{\rm CH_4} = A\,T^{C}\exp\!\left(\frac{B}{T}\right),
\end{equation}
\begin{equation}
  C = -\frac{M}{R}\,\Delta c_p, \qquad
  B = -\frac{M}{R}\left(L + \Delta c_p\,T_{\rm tp}\right), \qquad
  A = p_{\rm tp}\,T_{\rm tp}^{-C}\exp\!\left(-\frac{B}{T_{\rm tp}}\right),
\end{equation}
with $M/R = 16.043/8.3143$. The other constants depend on the phase:
\begin{itemize}
  \item solid ($T \le T_{\rm tp}$): $\Delta c_p = 2.213 - 2.650$ and latent
        heat of sublimation $L = 611.10$;
  \item liquid ($T > T_{\rm tp}$): $\Delta c_p = 2.213 - 3.370$ and latent
        heat of vaporisation $L = 552.36$;
\end{itemize}
with heat capacities in J\,g$^{-1}$\,K$^{-1}$ and latent heats in
J\,g$^{-1}$. The result is in bar, and the code returns dyn\,cm$^{-2}$.

\begin{thebibliography}{99}

\bibitem[Ackerman \& Marley(2001)]{ackerman2001}
Ackerman, A.~S., \& Marley, M.~S. 2001, ApJ, 556, 872,
\emph{Precipitating Condensation Clouds in Substellar Atmospheres}

\bibitem[Barshay \& Lewis(1976)]{barshay1976}
Barshay, S.~S., \& Lewis, J.~S. 1976, ARA\&A, 14, 81,
\emph{Chemistry of Primitive Solar Material}

\bibitem[Batalha et al.(2025)]{batalha2025}
Batalha, N.~E., Rooney, C.~M., Visscher, C., et al. 2025,
arXiv:2508.15102,
\emph{Condensation Clouds in Substellar Atmospheres with Virga};
code: \url{https://github.com/natashabatalha/virga}

\bibitem[Buck(1981)]{buck1981}
Buck, A.~L. 1981, J. Appl. Meteorol., 20, 1527,
\emph{New Equations for Computing Vapor Pressure and Enhancement Factor}

\bibitem[Charnay et al.(2018)]{charnay2018}
Charnay, B., B\'ezard, B., Baudino, J.-L., et al. 2018, ApJ, 854, 172,
\emph{A Self-consistent Cloud Model for Brown Dwarfs and Young Giant
Exoplanets: Comparison with Photometric and Spectroscopic Observations}

\bibitem[Ferguson \& Nuth(2008)]{ferguson2008}
Ferguson, F.~T., \& Nuth, J.~A. 2008, J. Chem. Eng. Data, 53, 2824,
\emph{Vapor Pressure of Silicon Monoxide}

\bibitem[Flatau et al.(1992)]{flatau1992}
Flatau, P.~J., Walko, R.~L., \& Cotton, W.~R. 1992, J. Appl. Meteorol.,
31, 1507, \emph{Polynomial Fits to Saturation Vapor Pressure}

\bibitem[Grant et al.(2023)]{grant2023}
Grant, D., Lewis, N.~K., Wakeford, H.~R., et al. 2023, ApJL, 956, L29,
\emph{JWST-TST DREAMS: Quartz Clouds in the Atmosphere of WASP-17b}

\bibitem[Lodders(1999)]{lodders1999}
Lodders, K. 1999, ApJ, 519, 793,
\emph{Alkali Element Chemistry in Cool Dwarf Atmospheres}

\bibitem[Lodders(2002)]{lodders2002}
Lodders, K. 2002, ApJ, 577, 974,
\emph{Titanium and Vanadium Chemistry in Low-Mass Dwarf Stars}

\bibitem[Lodders \& Fegley(1998)]{lodders1998}
Lodders, K., \& Fegley, B. 1998,
\emph{The Planetary Scientist's Companion} (New York: Oxford University Press)

\bibitem[Marley et al.(2021)]{marley2021}
Marley, M.~S., Saumon, D., Visscher, C., et al. 2021, ApJ, 920, 85,
\emph{The Sonora Brown Dwarf Atmosphere and Evolution Models. I.}

\bibitem[Morley et al.(2012)]{morley2012}
Morley, C.~V., Fortney, J.~J., Marley, M.~S., et al. 2012, ApJ, 756, 172,
\emph{Neglected Clouds in T and Y Dwarf Atmospheres}

\bibitem[Morley et al.(2024)]{morley2024}
Morley, C.~V., Mukherjee, S., Marley, M.~S., et al. 2024, ApJ, 975, 59,
\emph{The Sonora Substellar Atmosphere Models. III. Diamondback}

\bibitem[Visscher et al.(2006)]{visscher2006}
Visscher, C., Lodders, K., \& Fegley, B. 2006, ApJ, 648, 1181,
\emph{Atmospheric Chemistry in Giant Planets, Brown Dwarfs, and Low-Mass
Dwarf Stars. II. Sulfur and Phosphorus}

\bibitem[Visscher et al.(2010)]{visscher2010}
Visscher, C., Lodders, K., \& Fegley, B. 2010, ApJ, 716, 1060,
\emph{Atmospheric Chemistry in Giant Planets, Brown Dwarfs, and Low-Mass
Dwarf Stars. III. Iron, Magnesium, and Silicon}

\bibitem[Wakeford et al.(2017)]{wakeford2017}
Wakeford, H.~R., Visscher, C., Lewis, N.~K., et al. 2017, MNRAS, 464, 4247,
\emph{High-temperature Condensate Clouds in Super-hot Jupiter Atmospheres}

\bibitem[Wetzel et al.(2013)]{wetzel2013}
Wetzel, S., Klevenz, M., Gail, H.-P., Pucci, A., \& Trieloff, M. 2013,
A\&A, 553, A92, \emph{Laboratory Measurement of Optical Constants of Solid
SiO and Application to Circumstellar Dust}

\bibitem[Wexler(1976)]{wexler1976}
Wexler, A. 1976, J. Res. Natl. Bur. Stand., 80A, 775,
\emph{Vapor Pressure Formulation for Water in Range 0 to 100\,$^\circ$C.
A Revision}

\bibitem[Wexler(1977)]{wexler1977}
Wexler, A. 1977, J. Res. Natl. Bur. Stand., 81A, 5,
\emph{Vapor Pressure Formulation for Ice}

\end{thebibliography}

\end{document}
